% Nature profile: one verified 2021 review, neutral source-inspired adaptation.
\PassOptionsToPackage{no-math}{fontspec}
\documentclass[UTF8,fontset=none,twocolumn,scheme=plain]{ctexart}
\usepackage[paperwidth=595.276bp,paperheight=790.866bp,left=40bp,right=37.276bp,top=49bp,bottom=42.866bp,headheight=20bp,headsep=4bp,footskip=23bp]{geometry}
% Shared dependencies and neutral material. Compile with XeLaTeX.
% No publisher logos, paper text, data, figures, or download metadata are included.
\usepackage{amsmath,amssymb,graphicx,xcolor,tikz,fancyhdr,needspace,xurl,hyperref,balance,etoolbox}
% The legacy balance package omits the first split box's final baseline depth.
% Reserve that measured depth rather than hiding a CJK overfull-box warning.
\makeatletter
\patchcmd{\@BAlancecol}{\setbox\@leftcolumn=\vbox to \@tempdima}{\advance\@tempdima by \dp1\relax\setbox\@leftcolumn=\vbox to \@tempdima}{}{\PackageWarning{cns-template}{Balance depth patch unavailable; inspect final columns}}
\makeatother
\usetikzlibrary{arrows.meta,positioning}
\usepackage[protrusion=true,expansion=false]{microtype}
\IfFontExistsTF{Noto Serif CJK SC}{%
  \setCJKmainfont{Noto Serif CJK SC}%
}{%
  \setCJKmainfont[BoldFont=FandolSong-Bold.otf,ItalicFont=FandolKai-Regular.otf]{FandolSong-Regular.otf}%
}
\IfFontExistsTF{Noto Sans CJK SC}{%
  \setCJKsansfont{Noto Sans CJK SC}%
}{%
  \setCJKsansfont[BoldFont=FandolHei-Bold.otf]{FandolHei-Regular.otf}%
}
\setmainfont{TeX Gyre Termes}\setsansfont{TeX Gyre Heros}
\setmonofont{TeX Gyre Cursor}
\definecolor{JournalBlue}{HTML}{0097C9}
\definecolor{JournalRed}{HTML}{B83443}
\definecolor{ScienceRed}{HTML}{CD1021}
\definecolor{CiteBlue}{HTML}{0000EE}
\hypersetup{colorlinks=true,linkcolor=JournalBlue,urlcolor=JournalBlue,
  pdftitle={CNS论文汉化：中性版式示例},pdfauthor={示例模板；非期刊官方模板},
  pdfsubject={Neutral source-inspired Chinese layout demonstration; no real research claims}}
\urlstyle{same}\raggedbottom
\setlength{\parindent}{1em}\setlength{\parskip}{0pt}
\setlength{\emergencystretch}{2em}\tolerance=2200
\XeTeXlinebreaklocale "zh"\XeTeXlinebreakskip=0pt plus .1pt
\pagestyle{fancy}\fancyhf{}
\renewcommand{\headrulewidth}{0pt}\renewcommand{\footrulewidth}{0pt}
\newcommand{\PaperTitle}{中文题名占位：从组织结构到证据解释}
\newcommand{\SampleAuthors}{Author A\textsuperscript{1}, Author B\textsuperscript{2,*}}
\newcommand{\SampleAffiliations}{\textsuperscript{1}作者单位一（示例占位）。\quad
  \textsuperscript{2}作者单位二（示例占位）。\par
  *通信作者信息待填。DOI、收稿和出版日期均应从目标原文核实。}
\newcommand{\SampleAbstract}{本段是中性摘要占位文字，不陈述任何真实研究发现。它演示中文题名、作者信息、摘要和双栏正文之间的层级关系。实际使用时，应完整替换为目标论文的忠实中文译文，并保留研究对象、方法、结果和结论的限定条件。图注、参考文献、作者声明和来源信息均属于独立的覆盖项。模板中的图形仅说明排版结构，不表示生物学机制、实验结果或统计关系。}
\newcommand{\DemoCaption}{\textbf{图1. 中性证据链示意（占位图，无实验数据）。}
  \textbf{A}，输入节点表示需要核对的原文材料。
  \textbf{B}，处理节点表示逐段转换及复核。
  \textbf{C}，输出节点表示正文、图注和书目等独立交付部分。
  箭头仅表示本示例的阅读顺序，不代表因果关系、效应方向或数量变化。
  该图由模板自行绘制，所有节点均为文字占位；无样本量、误差线、统计检验或比例尺。
  实际论文图注应逐项保留原稿报告的样本量、单位、检验方法、限定及图形署名。}
\newcommand{\DemoVector}{%
 \resizebox{\linewidth}{!}{%
 \begin{tikzpicture}[font=\sffamily\small,node distance=18mm,
   box/.style={draw=black!70,rounded corners=2pt,fill=black!3,minimum width=26mm,minimum height=16mm,align=center},
   flow/.style={-{Stealth[length=2.5mm]},line width=.7pt}]
  \node[box] (a) {\textbf{A}\quad 输入\\原文材料};
  \node[box,right=of a] (b) {\textbf{B}\quad 处理\\转换与复核};
  \node[box,right=of b] (c) {\textbf{C}\quad 输出\\中文译稿};
  \draw[flow] (a)--(b);\draw[flow] (b)--(c);
  \node[below=11mm of a,align=center] {标题与摘要\\正文与图表};
  \node[below=11mm of b,align=center] {保留限定\\检查单位};
  \node[below=11mm of c,align=center] {可追溯引文\\明确原文信息};
  \node[below=29mm of b,font=\sffamily\footnotesize,text=black!65] {示意占位；不包含真实实验数据};
 \end{tikzpicture}}%
}
\newcommand{\SampleBibliography}{%
 \section*{参考文献（示例占位）}
 {\fontsize{7.6bp}{10bp}\selectfont\setlength{\parindent}{0pt}\interlinepenalty=10000
 \hypertarget{ref:1}{}1. Author A. (20XX). Example reference title. \textit{Journal placeholder}. 卷、页码与DOI待填。此条仅供演示引文跳转。\par\smallskip
 \hypertarget{ref:2}{}2. Author B. (20YY). Another example reference title. \textit{Journal placeholder}. 出版信息待填；不要把占位条目带入正式译稿。\par}
}
\newcommand{\SampleSourceInfo}{%
 \section*{原文信息}
 以下字段全部为示例占位：作者；英文原标题；期刊；年份；卷期；页码或文章号；DOI或来源链接。实际使用时必须逐项替换并核对。\par
 本文件是\JournalName 风格的非官方中文译稿样板，依据一次已核验文章版式进行结构化改编，不代表该期刊所有年代或文体的统一标准，也不保证像素级复刻。无真实论文内容或期刊标识图像。\par
}
% Deliberately neutral demo prose, not translated from a published paper.
\newcommand{\SampleOpening}{%
 \section*{引言：示例正文与引文}
 这是一段用于检查双栏阅读节奏的中性示例文字。段落中的术语、缩写和标点需要在实际项目中与源文保持一致。中文译稿应呈现原论文的完整论述，而不是把多段证据压缩成一段摘要。这里的引文只链接到本文件末尾明确标为占位的参考文献\SampleCiteOne。\par
 若原句仅表示相关、可能或提示，译文也应保留相同的证据强度。数字、单位、正负号、上下标和范围端点应逐项核对。原文存在内部矛盾时，应分别保留相应位置的陈述，再用简短译注说明；不能根据常识把不同的数字静默统一。\par
 \subsection*{内容层级与跨页衔接}
 此小标题演示二级层级。实际论文可能采用不同的小标题密度，模板不会替代对原文结构的识别。双栏文章需按左栏、右栏和跨页续句的顺序读取。若段落在页末中断，应在完成整句核对后接入下一页，避免把图中文字误当正文，或重复翻译同一段。\par
 图表具有独立的信息范围。正文引用到某幅图，并不意味着其图注可以省略。图注中的每个分面、比较对象、时间点、样本量及统计说明都应对应源文。英文图内标签可以保留，并在图旁或文末提供中文释义；实际采用哪种安排，需要依据图形尺寸和译文长度判断\SampleCiteTwo。\par
 \subsection*{来源与可维护数据}
 本段继续演示正文密度。可编辑文本和图注数据应与最终版面分开维护，以便修订时准确定位变化。若只调整版式，应能证明正文、图注和书目内容没有被意外改写。更换字体或纸张后，须重新检查断行、图注与图片的相对位置以及页脚可读性。\par
 文内引用与书目之间应有明确对应关系。这里的两条演示书目没有真实DOI，也没有虚构的实验信息。使用模板时可增删条目，但必须同时更新编号、链接目标和引用位置。页码只表示中文译稿自己的分页，不用于冒充原出版页码。\par
 \subsection*{结果叙述的演示占位}
 本节的文字仅用于呈现长段落在双栏页面中的流动，不报告真实结果。实际翻译应保留观察对象、实验条件和结局指标之间的关系。例如，原文分别描述体外实验与动物实验时，不能把两类结果合并为同一层面的结论；研究使用的对照、时间顺序和分析范围也必须保留。图表中的比较范围不应扩大为对整个人群或全部疾病阶段的判断。\par
 论述中的转折与限制同样属于正文。若研究先提出一种解释，再讨论另一种可能，译文应沿用这一推进顺序。作者明确写出的不确定性应能在中文中被识别；如果尚未确定某个现象的方向或机制，就不能改成已经建立的因果关系。对于术语近似但研究层级不同的概念，可在首次出现时保留缩写，后文保持统一。\par
 \subsection*{方法信息与可复核性}
 方法部分往往包含密集数字和专有名词。翻译时须区分浓度、总量、频率和处理时长，不能因为中文句式调整而改变参数的归属。试剂名称、软件版本和数据访问标识通常保留原样；上下限、排除规则和显著性标准则需要逐字符核对。这里没有提供任何真实实验参数，避免把模板占位误认为可以直接执行的方法。\par
 表格与补充材料的引用也需要保留原层级。正文对补充图的指向并不意味着模板中已经附有该补充图；若源文件只包含主文，应如实标注所交付的范围。图表编号、附录顺序和参考文献数量都应由目标文章决定，本示例中的结构不构成固定配额，也不要求所有论文采用同样的章节。\par
}
\newcommand{\SampleClosing}{%
 \section*{讨论：版式适配与范围}
 上方矢量图是模板原创的排版占位图，不来自任何论文。图形与完整演示图注被放在同一页，便于读者直接核对分面说明。真实项目应使用经过授权或用户提供的原图素材，并保留图注署名与许可信息。尺寸不足时，可适度缩放或调整栏宽，不以删除图注解决空间问题。\par
 模板仅记录一次已核验文章的可观察设计特征。即使文章出自同一期刊，不同年代、研究论文与综述、校正清样与正式版本也可能采用不同的页眉和引文形式。应以用户提供的实际PDF为准，不能仅凭期刊名称推断配色、字号或引用方式\SampleCiteOne。\par
 \subsection*{演示检查项}
 本样板用于展示可编译性、中文字形、标题层级、跨栏图注、页眉页脚和内部链接。图中没有实际数据，因此不应被当作科学插图引用。替换内容后，还需检查所有页的渲染结果，确认无截字、碰撞、异常空白和遗漏段落。\par
 \SampleBibliography
 \SampleSourceInfo
}


\usepackage{eso-pic}
\newcommand{\JournalName}{Nature}
\hypersetup{hidelinks}
\setlength{\columnsep}{12bp}\setlength{\parindent}{0pt}
\renewcommand{\normalsize}{\fontsize{9.05bp}{11.95bp}\selectfont}
\ctexset{section={format=\bfseries\fontsize{10.5bp}{13bp}\selectfont,beforeskip=10bp,afterskip=3bp},subsection={format=\bfseries,beforeskip=7bp,afterskip=2bp}}
\newcommand{\SampleCiteOne}{\textsuperscript{\hyperlink{ref:1}{1}}}
\newcommand{\SampleCiteTwo}{\textsuperscript{\hyperlink{ref:2}{2}}}
\fancyhead[L]{\bfseries\fontsize{14bp}{18bp}\selectfont Review}
\renewcommand{\headrulewidth}{.3bp}
\fancyfoot[L]{\sffamily\fontsize{7.5bp}{9bp}\selectfont\textbf{中文译页\ \thepage}\quad |\quad Nature}
\fancyfoot[R]{\sffamily\fontsize{6.8bp}{9bp}\selectfont 非官方中文译稿示例}
\AddToShipoutPictureBG{\AtPageLowerLeft{\ifnum\value{page}<3\put(34,43){\rule{.3bp}{719bp}}\else\put(33.9283,47.3507){\rule{.35bp}{715.169bp}}\fi}}
\begin{document}\normalsize
\twocolumn[{\noindent{\bfseries\fontsize{25bp}{31bp}\selectfont\PaperTitle\par}\vspace{16bp}\hrule height .3bp\vspace{6bp}
\noindent\begin{minipage}[t]{.32\textwidth}\vspace{0pt}\sffamily\fontsize{7.5bp}{10bp}\selectfont
DOI：待核实\par\vspace{4bp}\hrule height .3bp\vspace{4bp}
收稿日期：待填\par 接受日期：待填\par 在线发表：待填\par\vspace{8bp}
本页所有研究信息均为占位。\par 非官方中文译稿样板。
\end{minipage}\hfill\begin{minipage}[t]{.656\textwidth}\vspace{0pt}
{\sffamily\bfseries\fontsize{9bp}{12bp}\selectfont\SampleAuthors\par}\vspace{8bp}
{\fontsize{10.4bp}{14.4bp}\selectfont\SampleAbstract\par}\vspace{7bp}
{\sffamily\fontsize{7.4bp}{10bp}\selectfont\SampleAffiliations\par}
\end{minipage}\par\vspace{8bp}\hrule height .3bp\vspace{11bp}}]
\balance\SampleOpening
\clearpage
\twocolumn[{\begin{center}\begin{minipage}{.78\textwidth}\DemoVector\end{minipage}\end{center}\vspace{3bp}
\noindent{\sffamily\fontsize{7.7bp}{10.3bp}\selectfont\DemoCaption\par}\vspace{13bp}}]
\balance\SampleClosing

% Additional Article / advance-online demonstration. The Review pages above
% retain identical text and page pixels; all new material below is neutral prose.
\clearpage\nobalance
\newgeometry{left=39.685bp,right=34.016bp,top=49bp,bottom=42.866bp,headheight=20bp,headsep=4bp,footskip=23bp}
\onecolumn
\fancyhead[L]{\bfseries\fontsize{14bp}{18bp}\selectfont Article}
\fancyfoot[L]{\sffamily\fontsize{7.5bp}{9bp}\selectfont\textbf{中文译页\ \thepage}\quad |\quad Nature}
\newcommand{\ArticleCite}[1]{\textsuperscript{\hyperlink{article-ref:#1}{#1}}}
\noindent{\bfseries\fontsize{25bp}{31bp}\selectfont 中文研究题名占位：从完整资料的连续记录\\到研究内容的可核验表达\par}
\vspace{16bp}\hrule height .3bp\vspace{6bp}
\noindent\begin{minipage}[t]{166.3bp}\vspace{0pt}\sffamily\fontsize{7.8bp}{10.5bp}\selectfont
DOI：按所给原稿核实\par\vspace{4bp}\hrule height .3bp\vspace{4bp}
收稿日期：按原稿填写\par\vspace{4bp}\hrule height .3bp\vspace{4bp}
接受日期：按原稿填写\par\vspace{4bp}\hrule height .3bp\vspace{4bp}
在线发表：按页面核对\par\vspace{4bp}\hrule height .3bp\vspace{7bp}
提前在线版无卷期页码时，明确记录当前版本，不自行补造。\par\smallskip
可见页面与提取文字不同时，需对照图像；若仍无法确定，应说明来源差异。
\end{minipage}\hfill\begin{minipage}[t]{343.9bp}\vspace{0pt}
{\sffamily\bfseries\fontsize{8.7bp}{11.5bp}\selectfont
示例作者甲\textsuperscript{1,7}，示例作者乙\textsuperscript{1,7,*}，示例作者丙\textsuperscript{2,7}，示例作者丁\textsuperscript{3}，示例作者戊\textsuperscript{4}，示例作者己\textsuperscript{5}，示例作者庚\textsuperscript{6}与示例作者辛\textsuperscript{1,*}\par}
\vspace{11bp}\hrule height .3bp\vspace{6bp}
{\fontsize{10.2bp}{14bp}\selectfont\setlength{\parindent}{0pt}
本页展示研究论文与提前在线版本的首页结构。题名、作者、日期、摘要和单位均为原创中性占位，不报告真实发现。不同文体可以共享期刊的黑白层级与栏宽，但出版字段及实际内容仍须按所给文章独立核对。正式译稿应完整保留研究问题、分析方法、结果及其限定，不能为了适配本例篇幅缩成摘要。\par\smallskip
日期与出处应具有明确依据。文件名、接受日期、PDF生成日期和DOI中的年份，均不能直接代替在线发表时间；未分配正式卷期页码时，也不应按另一篇文章补写。页面显示日期与底层提取文字不一致时，需要检查可见页面及其版本，记录判断所依据的信息。\par\smallskip
作者单位在本页以通栏脚注示意。实际机构名称、单位序号、共同贡献和通信作者对应关系，均应保持完整。前两页的综述演示与本页研究论文演示属于同一期刊的文体差异，不在总览中重复计为新期刊。\par\smallskip
正文、方法、扩展图表和报告摘要有各自的覆盖范围。只有实际收到的材料才能计为已翻译；正文中引到的独立附件若未提供，应如实说明。图内英文保留时，完整图注及标签释义仍应译成中文。此页全部内容只用于测试版面，不构成科研结果或实验建议。
\par}
\end{minipage}\par\vspace{10bp}\hrule height .3bp\vspace{9bp}
\noindent\begin{minipage}[t]{\dimexpr\linewidth/2-6bp\relax}\vspace{0pt}\normalsize
\section*{研究论文开篇（原创占位）}
实际论文常在摘要下接入双栏正文。译文应保留作者的论证顺序，不因标题或摘要长度变化而省去开篇段落。正文第一次实质讨论某幅图时，应能找到相应图像与完整图注\ArticleCite{1}。\par
本页短段落只演示栏结构，不预设完整文章的段数。科学原图中的坐标、单位、误差线、统计说明与颜色对应均需分别核对。
\par 科学叙述包含比较、否定和不确定性。原作者若把观察结果与可能解释分别陈述，译文也应保持这一层次，不把相关性改写成因果关系。段落在图表附近中断时，先恢复跨栏或跨页的完整句子，再决定译文分页。
\par 图注有独立的信息范围。即使图中已有刻度或统计符号，图注中的样本量、检验、区间、误差定义、缩写和署名也不能省略。保留科学原图时，可将中文标签释义放在文末，避免改写密集的数据图。
\par 此处全部文字为原创排版说明，不引用真实研究结果。替换为正式内容后，仍需逐段比对来源，并渲染检查所有页面；示例中的段数与长度并不是期刊限制。
\end{minipage}\hfill\begin{minipage}[t]{\dimexpr\linewidth/2-6bp\relax}\vspace{0pt}\normalsize
\section*{版本与内容的分别维护}
论文类型为Article时，页眉文字应随之改变；提前在线版与正式刊期版的出版身份也应分开记录。仅用同一期刊的既有模板，不能替代对目标版本的观察\ArticleCite{2}。\par
实际全文重新分页时，图文应按阅读顺序交替安排。不同纸型、字体与中文长度会改变页数，但不能改变研究数值、引文编号或声明的归属。
\par 出版信息与研究内容应分别维护。原文件可能包含已更新的可见日期，同时保留旧的底层文字；译稿需要说明所依据的页面版本。未报告的卷期、文章号或日期应保持缺失，不把文件生成时间当作发表时间。
\par 方法、数据与代码可用性以及版权许可等声明可能位于正式页码之后。只有逐页清点实际提供的内容，才能避免因主文结束而漏译这些部分。正文仅提到但没有随文件提供的附件，应明确列出范围限制。
\par 字体回退和中文行距属于版式适配。它们可以改变分页，但不应改变来源身份、科学内容或署名关系。任何局部样式调整之后，都要检查受影响页面，确认没有隐藏的截断和重叠。
\end{minipage}
\vfill\noindent\hrule height .3bp\vspace{5bp}
{\sffamily\fontsize{7.1bp}{9.2bp}\selectfont\noindent
\textsuperscript{1}作者单位一（原创中性占位）：正式稿填写完整机构与地点，不沿用本页示例。
\textsuperscript{2}作者单位二（中性占位）。\textsuperscript{3}作者单位三（中性占位）。
\textsuperscript{4}作者单位四（中性占位）。\textsuperscript{5}作者单位五（中性占位）。
\textsuperscript{6}作者单位六（中性占位）。\textsuperscript{7}共同贡献作者关系按来源填写。
*通信作者及联系方式按原文填写。本页无真实作者、机构、邮箱或研究数据。
仅供学术交流、出版物版权归原作者及出版商所有。\par}
\clearpage
\twocolumn[{\noindent\begin{center}
\begin{tikzpicture}[font=\sffamily\fontsize{9bp}{12bp}\selectfont,node distance=15bp,
 box/.style={draw=black!70,rounded corners=2bp,fill=black!3,minimum width=102bp,minimum height=46bp,align=center},
 flow/.style={-{Stealth[length=2.5mm]},line width=.7bp}]
\node[box] (a) {\textbf{A}\quad 主文\\论点与图表};
\node[box,right=of a] (b) {\textbf{B}\quad 方法\\定义与参数};
\node[box,right=of b] (c) {\textbf{C}\quad 扩展内容\\原始编号};
\node[box,right=of c] (d) {\textbf{D}\quad 报告摘要\\已填与未填};
\draw[flow] (a)--(b);\draw[flow] (b)--(c);\draw[flow] (c)--(d);
\end{tikzpicture}\end{center}\vspace{3bp}
{\sffamily\fontsize{7.6bp}{9.2bp}\selectfont\noindent
\textbf{图2｜材料层级与范围记录（原创中性占位，无实验数据）。}
\textbf{A}，主文包含完整论述、图像与图注。\textbf{B}，方法中的定义、参数与限定具有独立内容范围。
\textbf{C}，实际提供的扩展内容沿用原编号。\textbf{D}，报告摘要需区分作者已填答案、勾选状态和未填写的模板提示。
箭头只示意本例阅读顺序，不表示实验流程、因果关系或数量变化；本图不含真实结果、样本量、误差线或比例尺。\par}\vspace{12bp}}]
\balance
\section*{方法与扩展内容（中性示例）}
方法段落中的定义、软件版本、比较对象、阈值及单位需对应所给来源。具体步骤和数值的翻译不由模板生成，也不按另一篇论文补齐。正文与方法中同一术语如果作用范围不同，应保留各自限定\ArticleCite{1}。\par
数据与代码可用性、伦理、经费、作者贡献、利益冲突和出版许可等声明属于独立内容范围。未提供的材料应如实记录，不把未填写的表格提示改写成已实施的研究。\par
\subsection*{扩展图与报告摘要}
扩展图可能是原生矢量，也可能已经是栅格图。保留原图时须检查全部面板、边界、图例、尺度和署名；放大并不会使原始栅格图获得额外信息。长图注可跨相邻页面续排，但应标明续注并完整保留尾句。\par
报告摘要可作为研究者填写时的中文参考。原研究的回答、不适用、确认勾选和空白提示需分别标注；不得将本页中性示例直接作为任何研究的填写答案。
\subsection*{原图清理的复核}
真正删除下载来源需检查实际承载对象。只改变裁区或在页面上遮白，不能证明隐藏对象已经移除。处理后还应比较科学图的可见文字与像素，确认没有意外改变字符宽度、面板顺序或颜色。\par
若某种通用清理方法使原有矢量文字改变，应采用更精确的对象处理并重新核验，不以肉眼近似代替记录\ArticleCite{2}。
\subsection*{可编辑工程与检查记录}
正文、图注、书目、图像及标签释义应具有清楚的文件边界，便于后续修订。独立源包须保留重建所需的相对路径与依赖，避免将本机目录、旧版输出、原始全文或临时提取稿一并分发。\par
科学内容检查与版面检查应分别进行。逐段对照确认无漏译或重复，并不等于图像边界完整；同样，页面看起来整齐，也不能证明引文链接、隐藏资源或书目编号都正确。内容变化、裁图变化与字体变化应有可区分的记录。
\subsection*{长文分页的调整}
人工分页仅用于本例展示。实际文章需按题名、摘要、图表高度及完整图注自然重排，尽量将图像放在相应论述附近；不能为了保留原页数而缩小到不可读字号。编译、渲染和源包独立重建通过后，才形成最终交付版本。
\section*{声明（原创占位）}
本样板没有实际研究、患者、实验对象、资助或作者声明。实际翻译按原文完整保留各项声明及其限定。本示例的图文和表述均为原创中性内容，不代表期刊官方模板或出版认可。
\section*{研究论文示例书目}
{\fontsize{7.6bp}{9.6bp}\selectfont
\noindent\hypertarget{article-ref:1}{1.} Author Placeholder. Source coverage and traceable content. Original neutral reference placeholder; no real publication.\par
\noindent\hypertarget{article-ref:2}{2.} Editorial Placeholder. Version-aware figure and layout checking. Original neutral reference placeholder; no DOI assigned.\par}
\section*{原文信息（研究论文占位）}
作者、英文原题、在线发表日期、卷期、原刊页码与DOI均应按目标来源填写。当前为原创中性排版示例，不对应真实论文；提前在线版尚无的正式出版字段不编造。\par
仅供学术交流、出版物版权归原作者及出版商所有。此页与前述综述演示分开示意内容范围，四页合计仍为一套Nature期刊样板。

\end{document}
